Lembremos que pela definição \ref{def:grupo-completo} temos que um grupo $G$ é dito completo  se $Z(G)=1$ e $Aut(G)=Inn(G)$. Assim, temos que se $G$ é completo, então temos que todo automorfismo de $G$ é interno e pelo corolario \ref{cor:1.2.4} temos que $G\cong Aut(G)=Inn(G)$.
Portanto, nosso interesse será verificar se $S_n$ com $n\geq 3$ e $n\neq 6$ é um grupo completo.\\

Para isso, vamos relembrar alguns dos resultados e propriedades importantes que precisaremos para provar isso.
\section{Propriedades do $S_n$}
Considere $n\in\mathbb{N}$, $n\geq 3$, então lembremos
\begin{align*}
  S_n=\{\sigma:\{1,\cdots,n\}\to\{1,\cdots,n\}:\sigma\text{ é bijetiva.}\}
\end{align*}
Sabemos que $|S_n|=n!$.\\
Então, dado  $\sigma\in S_n$ com $\sigma\neq 1$, podemos considerar a estructura cíclica de $\sigma$.
\begin{definition}
  Chamamos a $\sigma$ um cíclo de comprimendo $m$ se $\sigma$ é uma permutação do tipo
  \begin{align*}
    \sigma:G&\to G\\
    i_1&\to i_2,\\
    i_2&\to i_3,\\
    &\vdots\\
    i_m&\to i_1.
  \end{align*}
  Escrevemos $\sigma=(i_1,i_2,\cdots,i_m)$ ou $m$-cíclo. 
\end{definition}
Note que toda permutação $\sigma\neq 1$ em $S_n$ pode ser escrita como um produto de cíclos disjuntos (cíclos que não mvem elementos em común).\\
Dois cíclos disjuntos comutam entre sí, então ela é a estructura cíclica de $\sigma$, i.e, é a manera única a mas de ordem do $\sigma$ onde $\sigma$ se escreve como produto de cíclos disjuntos.
Exemplos: 
\begin{align*}
  (9357)=(3579)=(5793)=(7935)\\
  (3\,5\,7\,9)(2\,4)(8\, 10\, 11)=(3\, 5\, 7\, 9)(9\, 10\, 11)(2\, 4)
\end{align*}
Você pode fazer e seguinte exercicio.
\begin{align*}
  (359)(89)(34)=(34598).
\end{align*}
\begin{note}
  Todo $m$-cíclo pode ser escrito como um produto de [transposições] (cíclos de comprimendo $2$) não necesariamente disjuntos, por exemplo
  \begin{align*}
    (2359)=(29)(25)(23).
  \end{align*}
  Geralmente 
  \begin{align*}
    (i_1\, i_2\, \cdots\, i_m)=(i_1i_m)(i_1i_{m-1})\cdots (i_1i_2).
  \end{align*}
\end{note}
\begin{note}[Propriedades de $S_n$]\label{note:propriedades-gerador-Sn}
  Um conjunto gerador de $S_n$
  \begin{align*}
    \left\langle (ij):1\leq i\neq j\leq n \right\rangle=S_n.
  \end{align*}
  Temos as seguintes propriedades
  \begin{itemize}
    \item $S_n=\left\langle (12),(23),\cdots,((n-1) n) \right\rangle$.
    \item $S_n=\left\langle (12),(12\cdots n) \right\rangle$.
    \item Dada $\varsigma=(ij)$ transpoção com $i\leq j-1$, temos $(ij)=((j-1)j)(i(j-1))((j-1)j)$. 
  \end{itemize}
\end{note}
\begin{definition}[Permutações pares e impares]
  $\sigma\in S_n$ é \textbf{par} se pode ser escrita como um produto de uma cantidade par de transposições. Caso contrário, $\sigma$ é \textbf{impar}. 
\end{definition}
\begin{definition}[$A_n$]
  Definimos $A_n$ como
  \begin{align*}
    A_n:=\{\sigma\in S_n:\sigma\text{ é par.}\}
  \end{align*}
  Ele é um subgrupo de $S_n$ dito subgrupo \textbf{alternado de grau $n$} de $S_n$.\\
  Note que $|A_n|=\frac{n!}{2}$.
\end{definition}
\begin{note}[$A_n\normal S_n$]
  Note que $A_n\normal S_n$ para todo $n\geq 5$.\\
  Seja $\varsigma=(ij)$ uma transposição de $S_n$ defina
  \begin{align*}
    f:S_n&\to A_n\\
    \sigma&\to f(\sigma)= 
    \begin{cases}
      \sigma, &\text{ se } \sigma\text{ e par} \text{,} \\
      \varsigma\sigma, &\text{ se } \sigma\text{ e impar} .
    \end{cases}
  \end{align*}
  Logo $f$ é um homomorfismoo sobrejetivo, $\ker f=:\{1,\varsigma\}=\left\langle \varsigma \right\rangle$. Assim, $|S_n/\ker f|=2\cong \Img f$, logo $[S_n:A_n]=2$ e portanto $A_n\normal S_n$. 
\end{note}
Consequentemente $S_n$ não é um grupo simple, mas, vamos mostrar que $A_n$ é um grupo sumple com $n\geq 3$ e $n\neq 4$.\\ 
Em $S_4$ nos temos que $|A_4|=12$, logo $\tilde{K_4}=\{1,(12)(34),(13)(23),(14)(23)\}$ e $\tilde{K_4}\normal A_4$ (mais do fue isto $\tilde{K_4}\normal S_4$).
\begin{proposition}
  Duas permutações $\alpha,\beta\in S_n$ das conjugada se, e somente se, $\alpha$ e $\beta$ tem a mesma estructura cíclica.
\end{proposition}
\begin{proof} 
  Suponha que existe $\gamma\in S_n$ tal que $\alpha=\beta^{\sigma}$, vamos admitir que $\alpha$ é um $m$-cíclo, digamos $\alpha=(i_1i_2\cdots i_m)$, temos
  \begin{align*}
    \beta\gamma(i_1)&=\gamma\alpha(i_1)\\
    &=\gamma(i_2).\\
    \beta\gamma(i_2)&=\gamma\alpha(i_2)\\
    &=\gamma(i_3).\\
    &\vdots
  \end{align*}
  Concluimos que $\beta$ é o cíclo $(\gamma(i_1),\gamma(i_2)\cdots\gamma(i_m))$.

  Reciprocamente, sejan
  \begin{align*}
    \alpha&=(i_1^{(1)}i_2^{(1)}\cdots i_{k_1}^{(1)})\cdots(i_1^{(s)}i_{2}^{(s)}\cdots i_{k_s}^{(s)}),\text{ disjunto }\\
    \beta&=(j_1^{(1)}j_2^{(1)}\cdots j_{k_1}^{(1)})\cdots(j_1^{(s)}j_{2}^{(s)}\cdots j_{k_s}^{(s)}),\text{ disjunto }
  \end{align*}
  Defina $\gamma: i_t^{l}\to j_t^{l}$.
\end{proof}
\begin{theorem}[$A_n$ é simple]
  Seja $G=A_n$. Se $n\geq 5$ temos que $A_n$ é simple. 
\end{theorem}
\begin{proof} 
  Seja $\{1\}\neq N\normal A_n$, para mostrar que $N=A_n$ basta mostrar que $N$ contem um $3$-cíclo, ja que os $3$-cíclos são conjugação de outro $3$-cíclo.\\
  Considere $p$ primo tal que $p\mid |N|$.\\
  Pelo teorema de Cauchy, existe $\sigma\in N$ tal que $o(\sigma)=p$.\\
  Escrevemos $\sigma$ como produto de ciclos disjuntos $\varsigma_i$, $1\leq i\leq m$, logo $\sigma=\varsigma_1\cdots\varsigma_m$.\\
  Deta forma, cada $\varphi_i$ é um $p$-cíclo.\\
  Analisamos os valorespossivels para $p$.
  \begin{enumerate}
    \item $p=2$.\\
      uma transposiçõe e assim, $m\geq 2$ par, escreva
      \begin{align*}
        \sigma=(ab)(cd)\varsigma_3\cdots\varsigma_m\in N.
      \end{align*}
      Considere $\gamma=(abc)$, temo que como $N$ é normal se pode calcular
      \begin{align*}
        \sigma^{\gamma}\sigma^{-1}&=(ab)^{\gamma}(cd)^{\gamma}\varsigma_3^{\gamma}\cdots\varsigma_{m}^{\gamma}\cdot \sigma^{-1}\\
        &=(ab)^{\gamma}(cd)^{\gamma}(ab)(cd)\\
        &=(ac)(bd)\in N.
      \end{align*}
      Agora, como $n\geq 5$, podemos tomar $k\notin\{a,b,c,d\}$, consideremos $\theta=(ack)$ e $\alpha=(ac)(bd)$ e temos que $(akc)=\alpha^{\theta}\alpha\in N$.
    \item $p=3$.\\
      Neste caso, cada $\varsigma_i$ é um $3$-cíclo, se $m=1$, acabou. Caso contrário, se $m\geq 2$, escreva $\sigma=(abc)(def)\varsigma_3\cdots\varsigma_m$. Considere o $3$-cíclo $\gamma=(bcd)$ e, como $N$ é normal se cumpre que $\sigma^{\gamma}\sigma^{-1}\in N$ e é um $5$-cíclo, este caso será resolvido a partir de na próxima situação.
    \item $p>3$.\\
      Consider $\varsigma_1=(a_1a_2\cdots a_p)$: $p$-cíclo em $N$ e $\gamma=(a_2a_3a_4)$ (observe que $\gamma$ comuta com todos $\varsigma_i$ com $i\geq 2$). Temos que
      \begin{align*}
        \sigma^{\gamma}\sigma^{-1}=(a_2a_3a_5)\in N.
      \end{align*}
      Isto mostra o resultado.
  \end{enumerate}
  Em conclução $N=A_n$.
\end{proof}

\section{$S_n$ é completo com $n\geq 3$ e $n\neq 6$}
A seguir, temos um lema técnico que será útil para verificarmos que, quando $n\geq 3$ e $n\neq 6$, todos os automorfismos de $S_n$ são internos. 
\begin{lemma}\label{lema:2.2.3}
  Sejam $n\neq 6$ e $\tau\in S_n$ um elemento de ordem $2$ que é produto de $k$ transposições disjuntas. Então
  \begin{enumerate}
    \item O tamanho da classe de conjugação de $\tau$ é $\displaystyle \frac{n!}{2^{k}(n-2k)k!}$.
    \item Se $k>1$ então $|Cl(\tau)|\neq |Cl(\sigma)|$, para qualquer transposição $\sigma$ de $S_n$. 
  \end{enumerate}
\end{lemma}
\begin{proof} 
  Note que como $|S_n|=n!$ e $2\mid n!$, pelo teorema de Cauchy, temos que existe $\tau\in S_n$ que cumpre que $o(\tau)=2$. Assim, o grupo $S_n$ possui um elemento de ordem $2$ que é produto de $k$ transposições disjuntas com $1\leq k\leq\frac{n}{2}$, isto porque nos podemos ver que como $o(\tau)=2$, temos que se $\tau(i_l)=j_l$, então $\tau(j_l)=i_l$, pelo que você pode considerar a transposição $(i_l\,j_l)$. Por outro lado, note que se $i_l=j_l$ para algúm $l$, isto é que $\tau$ fixa $i_l$, pelo que nos vamos ignorar esse caso. Assim, note que como cada transposição $(i_l\,j_l)$ tem $2$ elementos $i_l$ e $j_l$ que cumpren que $i_l\neq j_l$, como $S_n$ são permutacões de $n$ elementos, temos que $2k\leq n$, assim temos que $1\leq k\leq\frac{n}{2}$. Em geral, temos que $\tau$ se pode escrever como 
  \begin{equation}\label{eq:tau-lema1}
    \tau=(i_1\,j_1)(i_2 j_2)\cdots(i_k,j_k).
  \end{equation}
  Agora, note que se $\tau$ não fixa $2k$ elementos, temos que $\tau$ fixa $n-2k$ elementos.\\
  Assim, contaremos o tamanho da classe de conjugação de $\tau$ em termos de $k$ e $n$. Analisando, novamente, que $\tau$ se escreve como produto de $k$ transposições disjuntas \eqref{eq:tau-lema1}, então que as posibilidades de escrever a primeira transposição são $\frac{n(n-1)}{2}$, isto ja que pra primeira transposição $(i_1\,j_1)$ temos que $i_1$ tem $n$ posibilidades, como $j_1\neq i_1$, temos que $j_1$ tem $n-1$ posibilidades e como a transposição $(i_1\,j_1)=(j_1,i_1)$ temos que as posibilidades de escrever a primeira transposição $(i_1\,j_1)$ são $\frac{n(n-1)}{2}$. Portanto, a próxima posibilidade de escolha para a segunda transposição, como restam $n-2$ elementos, então será $\frac{(n-2)(n-3)}{2}$.\\
  Prosseguindo a contagem dessa forma, segue que a possibilidade total de escolhas será dada pelo produto das posibilidades, mas, dentre as possibilidades mencionadas acima, porém, como a ordem das transposições não importa, dividiremos o resultado desse produto por $k!$, já que esse é o número de vezes que as mesmas $k$ transposições serão misturadas. Então teremos
  \begin{align*}
    |Cl(\tau)|&=\frac{1}{k!}\cdot \prod_{i=1}^{k}\frac{(n-2(i-1))(n-(2i-1))}{2}\\
    &=\frac{1}{k!}\cdot\prod_{i=0}^{k-1}\frac{(n-2i)(n-(2i+1))}{2}\\
    &=\frac{1}{k!}\cdot\frac{(n)(n-1)\cdots(n-2k-2)(n-2k-1)}{2^{k}}\\
    &=\frac{n!}{2^{k}k!(n-2k)!},
  \end{align*}
  maneiras de escrever o produto de $k$ transposições disjuntas. Assim concluimos o item \textbf{1.} afirmando que
  \begin{equation}\label{eq:clase-tau-lema1}
    |Cl(\tau)|=\frac{n!}{2^{k}k!(n-2k)!}.
  \end{equation}

  Agora, para provarmos o item \textbf{2.}, note que se $k=1$, então temos que $\tau$ é uma tranposição $(i_1\,j_1)$ e de \eqref{eq:clase-tau-lema1} temos que
  \begin{align*}
    |Cl(\tau)|=\frac{n!}{2(n-2)!}.
  \end{align*}
  Assim, queremos garantir que para $k>1$, o tamanho da cclasse de $\tau$ é diferente da classe de uma transposição, istó é
  \begin{align*}
    |Cl(\tau)|=\frac{n!}{2^{k}k!(n-2k)!}\neq\frac{n!}{2(n-2)!}.
  \end{align*}
  Que é o mesmo que provar
  \begin{equation}\label{eq:item2-lema1}
    2^{k}k!(n-2k)!\neq 2(n-2)!,\quad\text{ para todo }1<k\leq\frac{n}{2}.
  \end{equation}
  \begin{itemize}
    \item Se $k=2$, é fácil ver que $(4)(2)(n-4)\neq 2(n-2)!$ se $n\geq 2$.
    \item Agora, suponha que $k\geq 3$, pelo que temos que $\tau$ pode sé produto de $3$ transposições disjuntas, mas  istó implica que $n\geq 6$.\\
      Primeiro, note que se $n=6$ e $k=3$ temos que
      \begin{align*}
        2^{3}(3!)(6-2(3))!=(8)(9)(1)=2(6-2)!
      \end{align*}
      o que mostra a necessidade de nossa hipótese $n\neq 6$.\\
      Agora, como $n>6$ para provar \eqref{eq:item2-lema1} mostrarems que
      \begin{align*}
        2^{k}k!(n-2k)!<2(n-2)!
      \end{align*}
      que é o mesmo que provar
      \begin{align*}
        2^{k-1}k!(n-2k)!<(n-2)!.
      \end{align*}
      Note que pela definição de número binomial temos que
      \begin{align*}
        1\leq \binom{n-k}{k}=\frac{(n-k)!}{k!(n-2k)!}.
      \end{align*}
      Daqui segue que $k!(n-2k)!\leq (n-k)!$ o que implica que
      \begin{equation}\label{eq:item2-lema1-2}
        2^{k-1}k!(n-2k)!\leq 2^{k-1}(n-k)!.
      \end{equation}
      Agora, vamos ver que 
      \begin{equation}\label{eq:item2-lema1-3}
        2^{k-1}<(n-2)(n-3)\cdots(n-(k-1))=\frac{(n-2)!}{(n-k)!}.
      \end{equation}
      Para mostrar isto, observe que no lado direito dessa inequação temos $(n-2)-(n-k)=k-2$ fatores com o de menor valor sendo $(n-(k-1))$. Assim, lembre que como $1<k\leq \frac{n}{2}$, temos que $k-1\leq\frac{n-2}{2}$, logo temos que como $n>6$ se cumpre que 
      \begin{align*}
        n-(k-1)\geq n-\frac{n-2}{2}=\frac{n+2}{2}>4.
      \end{align*}
      Assim, temos que
      \begin{align*}
        2^{k-1}=\underbrace{2\cdots 2}_{k-1\text{ vezes}}<\underbrace{4\cdots 4}_{k-2\text{ vezes}}\leq (n-2)\cdots(n-(k-1))=\frac{(n-2)!}{(n-k)!}.
      \end{align*}
      Assim, usando na inequação \eqref{eq:item2-lema1-3} em na inequação \eqref{eq:item2-lema1-2} temos que
      \begin{align*}
        2^{k-1}k!(n-2k)!\leq 2^{k-1}(n-k)!\leq (n-2)!
      \end{align*}
      o que prova na desigualdade \eqref{eq:item2-lema1} e conclui o item \textbf{2}.
  \end{itemize}
\end{proof}
Agora, estamos prontos para provar que $S_n$, para $n\neq 6$ e $n\geq 3$, é um grupo completo.
\begin{proposition}\label{prop:2.2.4}
  Para $n\geq 3$, temos que $Z(S_n)=\{1\}$.
\end{proposition}
\begin{proof} 
  Considere $\alpha\in S_n$ com $\alpha\neq 1$. Vamos mostrar que sempre existe $\tau\in S_n$ tal que $\alpha\tau\neq\tau\alpha$ e desta forma, temos que $\alpha\neq Z(S_n)$. Observe que sendo $\alpha\neq 1$ existe $i\in\{1,\cdots,n\}$ tal que $\alpha(i)=j\neq i$. Como $n\geq 3$, existe $k\in\{1,\cdots,n\}$, onde $k\neq i$ e $k\neq j$. Assim, tome $\tau=(i\,k)$, depois temos que
  \begin{align*}
    \alpha\tau(i)=\alpha(k)\neq j\,\text{e}\,\tau\alpha(i)=\tau(j)=j.
  \end{align*}
  Portanto, $\alpha\tau\neq\tau\alpha$, pelo que podemos conclui que $Z(S_n)=\{1\}$ com $n\neq 3$. 
\end{proof}
\begin{theorem}\label{theo:2.2.5}
  Se $n\neq 6$, então todo automorfismo de $S_n$ é interno. 
\end{theorem}
\begin{proof} 
  Seja $\sigma\in Aut(S_n)$ e $\psi$ uma transposição em $S_n$. Note que um automorfismo leva conjuntos com $m$ elementos em conjuntos com $m$ elementos, ou seja, $\sigma$ leva $\frac{n!}{2(n-2)!}$ elementos da classe de $\psi$ em $\frac{n!}{2(n-2)!}$ elementos. Assim, pelo lema \ref{lema:2.2.3}, temos que como na classe de $\sigma(\psi)$ tem $\frac{n!}{2(n-2)!}$ elementos, então a classe de $\sigma(\psi)$ só possui transposições, portanto podemos afirmar que $\sigma$ leva transposições em transposições.\\
  Tome $c_i=(i\,i+1)$ para $1\leq i\leq n-1$ e seja $\psi_i=\sigma(c_i)$. Para $i=1$ escreva $\psi_1=(a_1\,a_2)$. Note que $c_1=(1\,2)$ e $c_2=(2\,3)$ não comutam, temos que $\psi_1$ e $\psi_2$ também não comutam, pois $\sigma$ é um isomorfismo, istó e
  \begin{align*}
    \psi_1\psi_2=\sigma(c_1)\sigma(c_2)=\sigma(c_1c_2)\neq\sigma(c_2c_1)=\sigma(c_2)\sigma(c_1)=\psi_2\psi_1.
  \end{align*}
  Podemos escrever então $\psi_2=(a_2\,a_3)$ com $a_3\neq a_1$. Análogamente, $\psi_3$ não comuta com $\psi_2$, pois $c_3$ não comuta com $c_2$. Assim $\psi_3$ deve possuir $a_2$ ou $a_3$. No entanto, uma vez que $c_1$ comuta com $c_3$, $\psi_1$ comuta com $\psi_3$. Sendo assim, $\psi_3=(a_3\,a_4)$ onde $a_4\notin\{a_1,a_2,a_3\}$. Da mesma forma, $\psi_4$ possui interseção com $\psi_3$ e é disjunto de $\psi_1$ e $\psi_2$. Assim sendo $\psi_4=(a_4\,a_5)$, com $a_5\notin\{a_1,a_2,a_3,a_4\}$. Continuando deste mdo podemoos escrever $\psi_i=(a_i\,a_{i+1})$ com $a_1,\cdots,a_n$ distintos e $a_i\in\{1,\cdots,n\}$.\\
  Seja $\gamma\in S_n$ uma permutação tal que $\gamma(i)=a_i$. Então $(c_i)^{\gamma}=\gamma^{-1}c_i\gamma=\gamma^{-1}(i\,i+1)\gamma=(a_{i}\,a_{i+1})=\psi_i$. Agora, se tomamos $\alpha$ o automorfismo interno induzido por $\gamma$, ou seja
  \begin{align*}
    \alpha:S_n&\to S_n,\\
            \eta&\to\eta^{\gamma},
  \end{align*}
  nos podemos mostrar que $\alpha^{-1}\sigma$ fixa cada $c_i$. Isto é importante porque $\{c_i\}$ gera $S_n$ (\ref{note:propriedades-gerador-Sn}). Se $\alpha^{-1}\sigma$ fixar todo $c_i$, teremos que $\alpha^{-1}\sigma$ fixará todo elemento de $S_n$, ou seja, $\alpha^{-1}\sigma=Id$.\\
  Então, sabemos que $\psi_i=\sigma(c_i)$, logo temos que $\alpha^{-1}\sigma(c_1)=\alpha^{-1}(\psi_i)$, depois como
  \begin{align*}
    \alpha^{-1}:S_n&\to S_n,\\
               \eta&\to\eta^{\gamma^{-1}},
  \end{align*}
  e $\psi_i=c_i^{\gamma}$ temos que
  \begin{align*}
    \alpha^{-1}(\sigma(c_i))=\alpha^{-1}(\psi_i)=\alpha^{-1}(c_i^{\gamma})=(c_i^{\gamma})^{\gamma^{-1}}=c_i.
  \end{align*}
  Assim $\alpha=\sigma$ é podemos conclui que todo automorfismo $\sigma\in Aut(S_n)$ é um automorfismo interno, i.é, $\sigma\in Inn(G)$, como queríamos provar. 
\end{proof}
\begin{theorem}
  $S_n$ é um grupo completo, para todo $n\geq 3$ e $n\neq 6$.
\end{theorem}
\begin{proof} 
  Usando a proposição \ref{prop:2.2.4} e o teorema \ref{theo:2.2.5} temos que $Z(S_n)=\{1\}$ e além mais $Aut(S_n)=Inn(S_n)$, logo se pode conclui que $S_n$ é um grupo completo. 
\end{proof}

\section{$S_6$ não é completo}
Na seguinte parte, nos vamos provar que para $n=6$, não temos $S_6\cong Aut(S_6)$. Ou seja, $S_6$ tem um automorfismoque não é interno. Para demostrar isso, é necessário discutir certascaracterísticas do grupo $S_6$ e do seu grupo de automorfismos, $Aut(S_6)$.\\

Nos vamos provar que $S_6$ possui um subgrupo, que denotaremos de $K$ que cumpre que $|K|=120$ e que $K$ não contém transposições. Para começarmos  provar a existência do subgrupo $K$ de $S_6$ a seguinte definição deve ser levada em consideração.

\begin{definition}[Subgrupo transitivo]
  Dado um conjunto $X$, dizemos que um subgrupo $S$ de $Sim(X)$ é um \textbf{subgrupo transitivo} se dados $x,y\in X$, temos que existe $\sigma\in S$ tal que $\sigma(x)=y$. 
\end{definition}
\begin{lemma}\label{lema:1-S6-triste}
  Seja $G$ um grupo e $H\leq G$. Definamos $X=\{xH:x\in G\}$ (o conjunto das classes laterais a esquerda de $H$) e $\varphi$ da seguinte maneira
  \begin{align*}
    \varphi:G&\to Sim(X),\\
            g&\to \varphi_g:X\to X,\\
  \phantom{g}&\phantom{\to \varphi_g:}xH\to gxH.
  \end{align*}
  Então $\Img \varphi$ é um subgrupo transitivo de $Sim(X)$.
\end{lemma}
\begin{proof} 
  Seja $x,y\in X$, isto é que $x=aH$ e $y=bH$ com $a,b\in G$. Agora, nos vamos tirar $g\in G$ com $g=ba^{-1}$, note que se definimos $\sigma=\varphi_{g}\in\Img\varphi$, assim temos que
  \begin{align*}
    \sigma(x)=\varphi_{g}(aH)=gaH=ba^{-1}aH=bH=y.
  \end{align*}
  Logo como $x,y\in X$ são arbitrarios e $\varphi$ é um homomorfismo, temos que $\Img\varphi$ é um subgrupo transitivo de $Sim(X)$. 
\end{proof}
\begin{lemma}\label{lema:2-S6-triste}
  Existe um subgrupo transitivo $K$ de $S_6$ com ordem $120$ que não contém transposições. 
\end{lemma}
\begin{proof} 
  Considere $G=S_5$, $P\in Syl_5(S_5)$ e $H=N_{S_5}(P)=\{x\in G:xP=Px\}$. Então, como $|S_5|=5!$, nos temos que pelos teoremas de Silow $n_5\mid (2\ast 3\ast 4)=24$ logo temos que $n_5\in\{1,2,3,6,8,12,24\}$, mas como $n_5\equiv 1(mod\,5)$, temos que $n_5\notin\{2,3,8,12,24\}$, pelo que só temos que $n_5\in\{1,6\}$.\\
  Más note que $|\left\langle (12345) \right\rangle|=5$, $|\left\langle (12)(12345)(12) \right\rangle|=|\left\langle (13452) \right\rangle|=5$ e $\left\langle (12345) \right\rangle\neq\left\langle (13452) \right\rangle$, pelo que podemos afirmar que $n_5=6$.\\
  Assim, temos que $[G:H]=[S_5:N_{S_5}(P)]=6$.\\
  Portanto, definindo $X=\{aH:a\in G\}$ temos que $|X|=6$ e $|Sim(X)|=6!$, pelo que sabemos que $Sim(X)\cong S_6$. Deste modo consideremos
  \begin{align*}
    \varphi:S_5&\to Sim(X),\\
            g&\to \varphi_g:X\to X,\\
  \phantom{g}&\phantom{\to \varphi_g:}xH\to gxH.
  \end{align*}  
  Note que $\varphi$ é homomorfismo injetor, ja que
  \begin{align*}
    \ker\varphi=\{g\in S_5:\varphi_{g}=Id\}=\{g\in S_5:aH=gaH,\text{ para todo }a\in S_5\}.
  \end{align*}
  Mas note que se $a=1$, então temos que $H=gH$, pelo que podemos afirmar que $g\in H$, assim $\ker\varphi\leq H$ e portanto temos que $[S_5:H]\leq[S_5:\ker\varphi]$.\\
  Logo, como temos que pelo teorema de isomorfismos $\ker\varphi\normal S_5$, então como $A_5$ é simple, temos que $\ker\varphi=\{1\}$ ou $\ker\varphi= A_5$. Mas note que como $[S_5:H]=6\leq [S_5:\ker\varphi]$, então a única possibilidade é que $\ker\varphi=\{1\}$, pelo contrario se $\ker\varphi=A_5$, temos que $[S_5:A_5]=2\not\geq 6$. Portanto, temos que $\varphi$ é um homomorfismo injetor.\\
  Assim, definamos $K=\Img\varphi$ e como $\varphi$ é injetivo, temos que $|K|=|\Img\varphi|=|S_5|=120$ e pelo teorema de isomorfismos temos que $K\leq S_6$ que pelo lema \ref{lema:1-S6-triste} temos que $K$ é um subgrupo transitivo de $S_6$.

  Agora, nos vamos provar que $K$ não possui transposições.\\
  Raciocinemos pela contradição.\\
  Note que como $\varphi$ é um isomorfismo de $S_5$ sobre $\Img\varphi=K$, então temos que existe $\alpha\in K$ tal que $o(\alpha)=5$. Sem perdida de generalidade suponha $\alpha=(12345)$. Então, se nos assumimos que existe $(ij)\in K$ transposição. Como $K$ é transitivo e $j,6\in\{1,2,3,4,5,6\}$, então existe $\beta\in K$ tal que $\beta(6)=j$ e $\beta^{-1}(j)=6$. Logo note que $\beta(ij)\beta^{-1}\in K$, ou seja, $(\beta(i)\,6)$ com $\beta(i)\neq 6$, ja que $i\neq j$. Conjugando $(\beta(i)\,6)$ por todas las potencias de $\alpha$ teremos que $(16),(26),(36),(46),(56)\in K$, logo temos que como $\left\langle (16),(26),(36),(46),(56) \right\rangle= S_6$, então $K=S_6$, absurdo, pois $S_5\cong K$.\\
  Assim, temos que existe $K\leq S_6$ com ordem $120$ e que não possui transposições. 
\end{proof}
\begin{theorem}\label{th:S-6pronto}
  Existe um automorfismo de $S_6$ que não é interno. 
\end{theorem}
\begin{proof} 
  De acordo com o lema \ref{lema:2-S6-triste}, temos que existe $K\leq S_6$ com ordem $120$ e que não possui transposições. Note que como $|K|=120$ e $|S_6|=6\ast 120$, então temos que $[S_6:K]=6$, assim, considere $X=\{\alpha_1K,\alpha_2K,\alpha_3K,\alpha_4K,\alpha_5K,\alpha_6K\}$ as classes a esquerda de $K$. Então, suponha
  \begin{align*}
    \sigma:S_6&\to Sim(X),\\
             g&\to \sigma_g:X\to X,\\
   \phantom{g}&\phantom{\to \sigma_g:}\alpha_iK\to g\alpha_iK.
  \end{align*}
  Utilizando um argumento semelhante ao da demonstração do lema \ref{lema:2-S6-triste}, mostramos que $\sigma$ é injetiva. Além disso, note que como $|Sim(X)|=|S_6|$, então $\sigma$ é sobrejetiva. Portanto, $\sigma\in Aut(S_6)$.\\
  
  Raciocinando pela contradição, suponha que $\sigma\in Inn(S_6)$. Sendo assim, $\sigma$ é uma conjugação e preserva estructura cíclica dos elementos de $S_6$. Portanto, temos que $\sigma[(12)]=\sigma_{(12)}$ é uma transposição, entao temos que
  \begin{align*}
    \sigma:S_6&\to S_6,\\
             (12)&\to \sigma_(12):X\to X,\\
      \phantom{g}&\phantom{\to \sigma_(12):}\alpha_iK\to (12)\alpha_iK.
  \end{align*}
  para cada $1\leq i\leq 6$. Depois, como $\sigma_{(12)}$ é uma transposição, então $(12)\in \alpha_jK$ para algúm $1\leq j\leq 6$, isto é que $(12)\alpha_jK=\alpha_jK$, logo podemos conclui que $\alpha_j^{-1}(12)\alpha_j\in K$, mas note que como $\alpha_j^{-1}(12)\alpha_j$ é uma conjugação de a transposição $(12)$, entao $\alpha_j^{-1}(12)\alpha_j$ é uma tranposição em $K$, contradição, pois $K$ não possui transposições. Assim, temos que $\sigma\notin Inn(S_6)$, i.é, que $S_6$ possui um automorfismo que não é interno.
\end{proof}
\begin{corollary}[$S_6$ não é completo]
  $S_6$ não é um grupo completo. 
\end{corollary}
\begin{proof} 
  Note que pelo teorema \ref{th:S-6pronto} temos que $S_6$ possui um autormorfismo que não e interno, logo $Inn(S_6)\neq Aut(S_6)$ o que nos permite conclui que $S_6$ não é completo. 
\end{proof}

